\begin{abstract}
The 2018--2019 US trade war arrived after decades of efforts to deepen global economic ties through trade agreements, raising the question of whether these institutionalized relationships make countries more resilient to sudden, unilateral tariff shocks. This paper studies whether countries with bilateral trade agreements with the United States experienced different effects of tariffs on prices and volumes than countries without such agreements. Extending the event-study framework of Amiti, Redding, and Weinstein (2020), I construct a triple-differences design using monthly customs data at the country-product level, restricting attention to the first three tariff waves (Section 201 and Section 232) which applied non-discriminatorily across countries. I find that import volumes from agreement countries rose following tariff implementation while volumes from non-agreement countries fell: a 25\% tariff is associated with an approximately 25\% rise in imports from agreement countries and an approximately 44\% decline in imports from non-agreement countries twelve months after implementation. This effect is concentrated in steel products, where the volume gain for agreement countries is large and persistent; non-steel agreement countries experience only a transient volume gain that reverses after roughly six months, indicating that the resilience effect operates differently across product categories. The price evidence is consistent with the volume findings but cannot be causally identified due to pre-trend violations within product subsamples. These patterns are consistent with a cost-buffer mechanism in which prior tariff and non-tariff barrier reductions under existing agreements leave agreement-country exporters better positioned to absorb tariff shocks, complemented by expectations-based channels through which importers rebalance toward partners with more predictable policy environments.
\end{abstract}